\documentclass[10pt,a4paper]{amsart}
\usepackage[margin=19mm]{geometry}
\usepackage{amsmath,amssymb,mathtools}
\usepackage[scr=rsfs,cal=euler]{mathalfa}
\usepackage{xfrac}
\usepackage[unicode,hidelinks,bookmarks=false]{hyperref}
\newcommand{\ie}{i.e.\ }
\newcommand{\cf}{cf.\ }
\newcommand{\resp}{respectively\ }
\newcommand{\Subsection}{\subsection}
\providecommand{\nobreakdash}{\nobreak\hbox{-}}
\setlength{\emergencystretch}{2em}
\multlinegap0pt
\usepackage{amssymb}
\usepackage{caption}
\newtheorem{prop}{Proposition}
\title{Hermitian-Singer Functional and Differential Codes}
\author{G\'abor Korchm\'aros, Federico Romaniello, Valentino Smaldore}
\date{Source: arXiv:2511.14345v2 (CC BY 4.0)}
\begin{document}
\maketitle

\noindent\emph{Source and licence.} Hermitian-Singer Functional and Differential Codes. Version arXiv:2511.14345v2 (CC BY 4.0), distributed by arXiv under the Creative Commons Attribution 4.0 International licence, http://creativecommons.org/licenses/by/4.0/, as recorded in the arXiv licence metadata for this exact version and shown on the abstract page https://arxiv.org/abs/2511.14345v2. Full licence notice: \texttt{licenses/arxiv-2511.14345-CC-BY-4.0.txt}.\par\smallskip

\noindent\textit{Editorial scope.} Counted unit: the article's prop pro02112025A (source0.tex lines 181--183). The article's complete original proof of it is delivered with it (source0.tex lines 184--196, one \texttt{proof} environment of 13 lines). No mathematics is written, repaired or re-derived: the statement and the proof are the article's own lines, in the article's own notation, and no retained line is substituted or re-flowed.  Retained uncounted as the source-local context the proof needs: lines 151--153; lines 155--158.  Nothing was omitted from any retained span, so the package carries no omission record.  No retained line contains a citation, so the wrapper carries no bibliography and no bibliography entry.  No retained line contains a figure environment, an image-inclusion command or drawing code, so no image is delivered and the package declares no asset dependency.  Editorial formatting only, in addition: an AMS \texttt{amsart} preamble that restates the article's own declarations and macros used by the retained lines, namely the environments \texttt{prop} on separate counters, together with the standard packages those lines need.  The retained environments print in this document as Proposition 1 in order of appearance, whereas the article prints the same items with the numbering of its own full document.  The complete native source of the article as distributed in the arXiv e-print for this exact version is bundled unchanged as \texttt{sources/189353/source0.tex}.
\begin{prop}
\label{pro30102025} $H_t$ is a Hermitian curve of $\Pi$.
\end{prop}

\begin{equation}
\label{eq30102025}
1+t\gamma^{q^3+1}+t^{q^2+1}\gamma^{q^2-q+1}=0,\,\, \gamma^{q^4+q^2+1}=1.
\end{equation}

\begin{prop}
\label{pro02112025A} The $\Gamma$-orbit through any point of $\Pi$ is shared by exactly $q+1$ Hermitian curves $H_t$.  
\end{prop}

\begin{proof} We use the computation carried out in the proof of Proposition \ref{pro30102025}. 
%We have already shown that $H_t$ is a Hermitian curve. 
For a $q^4+q^2+1$-th root $\gamma$ of unity, the point $P=(\gamma:\gamma^{q^2+1}:1)$ of $\Pi$ belongs to $H_t$ if and only
if (\ref{eq30102025}) holds. Let $\delta=\gamma^{q^2-q+1}$. Since $t^{q^3+q}=t^{q+1}$, then $H_t$ passes through $P$ if and only if $1+t^q\delta^{q^2+q}+t^{q+1}\delta^q=0$. By $\delta^{q^2+q+1}=1$, the latter equation is equivalent to  
$\delta+t^q+t^{q+1}\delta^{q+1}=0$. Replacing $t$ by $t^{-1}$ gives $\delta t^{q+1}+t+\delta^{q+1}=0$. Therefore, it is enough to show that each root of the polynomial $P(X)=\delta X^{q+1}+X+\delta^{q+1}$ in $\overline{\mathbb{F}}$ is a ($q^2+q+1$)-th root of unity. This claims follows from $\delta P(X)^qX=X^{q^2+q+1}\delta^{q+1}+\delta X^{q+1}+\delta^{q^2+q+1}X$ and $\delta^{q^2+q+1}=1$.       
%Therefore, $(t^q,t^{q+1})$ is a solution of the linear system of equations
%\begin{align*} 
%&\delta^{q^2+q}X+\delta^qY+1=0;\\
%&X+\delta^{q+1}Y+\delta=0. 
%\end{align*}
%The determinant of this system equals $\delta^q-1$ which
%From the proof of Proposition \ref{pro30102025}, among the $q^2+q+1$-th root of unities, exactly $q+1$ are also roots of the polynomial $X^{q+1}+X+1=0$, and hence of the polynomial $X^{q^2+1}+X+1=0$.
\end{proof}

\end{document}
